% Source: ACM_Templates.pdf pp. 40-66. https://cp-algorithms.com/data_structures/sparse-table.html
\subsection{ST 表：静态区间最值}
\label{ds:st}
\AlgorithmInfo
    {在不修改数组时查询区间最大值。}
    {预处理时间、空间 $O(n\log n)$；查询 $O(1)$。}
    {非空数组按 1 下标传入，查询为非空闭区间。}
\subsubsection*{原理}
令 $f_{i,j}$ 为 $[i,i+2^j-1]$ 的最大值，合并相邻两个长度为 $2^{j-1}$ 的区间得到
\[
 f_{i,j}=\max(f_{i,j-1},f_{i+2^{j-1},j-1}).
\]
查询时取 $k=\lfloor\log_2(r-l+1)\rfloor$，用 $[l,l+2^k-1]$ 与 $[r-2^k+1,r]$ 覆盖原区间。两段可以重叠，因为最大值满足重复合并不改变答案。

\subsubsection*{使用教程}
\InlineCode{init(a)} 根据 \InlineCode{a[1..n]} 预处理；\InlineCode{query(l,r)} 返回最大值。更换数组后重新 \InlineCode{init}，旧状态会重置。\InlineCode{lg} 用整数递推，避免浮点对数的边界问题。
\begin{lstlisting}[style=mtycpp,numbers=none]
vector<long long> a={0,3,-2,7,4};
ST st;
st.init(a);
long long ans=st.query(2,4);
\end{lstlisting}
\needspace{7\baselineskip}
\subsubsection*{模板代码}
\CodeFile{code/data-structures/st.cpp}
\subsubsection*{需要注意}
可将两处 \InlineCode{max} 一起换成 \InlineCode{min} 或 \InlineCode{gcd}；普通求和不满足重复合并性质，不能直接套用本查询式。
